In this section we construct two radial Schwartz functions $a,b:\R^8\to i\R$ such that
\begin{align}\mathcal{F}(a)&=a\label{eqn:a-fourier}\\
    \mathcal{F}(b)&=-b\label{eqn:b-fourier}
\end{align}
with double zeroes at all $\Lambda_8$-vectors of length greater than $\sqrt{2}$. Recall that each vector of $\Lambda_8$ has length $\sqrt{2n}$ for some $n\in\N_{\geq 0}$.
We define $a$ and $b$ so that their values are purely imaginary because this simplifies some of our computations.
We will show in Section \ref{sec: g} that an appropriate linear combination of functions $a$ and $b$ satisfies conditions \eqref{eqn:g1}--\eqref{eqn:g3}.

Both functions are defined as $x \mapsto f(\|x\|^2)$, where $f : \R \to \C$ is a sum of six contour integrals of (quasi)modular forms against the Gaussian $e^{\pi i r z}$. The programme of this section is as follows. Bounds on the relevant (quasi)modular forms (obtained from Lemma~\ref{lemma:mod-div-disc-bound}) yield rapid decay of the integrals and their derivatives on $[0, \infty)$; smoothness on a neighbourhood of $[0,\infty)$ follows by differentiation under the integral sign; Schwartzness of $a$ and $b$ is then obtained from the smooth-cutoff criterion of Corollary~\ref{cor:schwartz-of-nonneg-decay}. The eigenfunction properties \eqref{eqn:a-fourier} and \eqref{eqn:b-fourier} follow from the Fourier transform of the Gaussian (Lemma~\ref{lemma:Gaussian-Fourier}) and the transformation laws of the integrands, and the double zeroes are exhibited by deforming contours.

\subsection{The Integrands for $a$}

To define the function $a$, we consider the following functions:
\begin{definition}\label{def:phi4-phi2-phi0}
\uses{def:Ek,def:E2}\lean{φ₀, φ₂', φ₄'}\leanok
\begin{align}
    \phi_{-4} &:= \frac{E_4^2}{\Delta} \label{eqn: def phi4} \\
    \phi_{-2} &:= \frac{E_4(E_2 E_4 - E_6)}{\Delta} \label{eqn: def phi2} \\
    \phi_{0} &:= \frac{(E_2 E_4 - E_6)^2}{\Delta} \label{eqn: def phi0}
\end{align}
\end{definition}

These functions are holomorphic on $\h$: their numerators are built from the holomorphic functions $E_2$, $E_4$ and $E_6$, and $\Delta$ does not vanish on $\h$ (Corollary~\ref{cor:disc-nonvanishing}).

% \begin{lemma}\label{lemma: phi fourier4 phi fourier2 phi fourier0}
% \uses{def:phi4-phi2-phi0}
%     These functions have the Fourier expansions
% \begin{align}
%     \phi_{-4}(z)\,=\,&q^{-1} + 504 + 73764\, q + 2695040\, q^2 + 54755730\, q^3 + O(q^4)\label{eqn: phi fourier4}\\
%     \phi_{-2}(z)\,=\,&720 + 203040\, q + 9417600\, q^2 + 223473600\, q^3 + 3566782080\, q^4+O(q^5)\label{eqn: phi fourier2}\\
%     \phi_{0}(z)\,=\,&518400\, q + 31104000\, q^2 + 870912000\, q^3 + 15697152000\, q^4+O(q^5)\label{eqn: phi fourier0}
% \end{align}
% where $q=e^{2\pi i z}$.
% \end{lemma}

The function $\phi_0(z)$ is not modular; however, it satisfies the following transformation rules:

\begin{lemma}\label{lemma:phi0-transform}\uses{def:phi4-phi2-phi0, lemma:Ek-is-modular-form, lemma:E2-transform-S, lemma:disc-cuspform}\lean{φ₀_periodic, φ₀_S_transform}\leanok
We have
\begin{align}
    \phi_0(z + 1) &= \phi_0(z) \label{eqn:phi0-trans-T} \\
    \phi_0\left(-\frac{1}{z}\right) &= \phi_0(z)-\frac{12i}{\pi}\,\frac{1}{z}\,\phi_{-2}(z)-\frac{36}{\pi^2}\,\frac{1}{z^2}\,\phi_{-4}(z).\label{eqn:phi0-trans-S}
\end{align}
\end{lemma}
\begin{proof}\leanok
\eqref{eqn:phi0-trans-T} easily follows from periodicity of Eisenstein series and $\Delta(z)$.
For \eqref{eqn:phi0-trans-S}, we use the weight $k$ transformation laws of $E_4$, $E_6$, $\Delta$ (of weights $4$, $6$, $12$) and the quasimodular transformation law of $E_2$ (Lemma~\ref{lemma:E2-transform-S}):
\begin{align}
    \phi_{0}\left(-\frac{1}{z}\right) &= \frac{(E_2(-1/z) E_4(-1/z) - E_6(-1/z))^{2}}{\Delta(-1/z)} \\
    &= \frac{((z^2 E_2(z) - 6iz / \pi) \cdot z^4 E_4(z) - z^6 E_6(z))^{2}}{z^{12} \Delta(z)} \\
    &= \frac{\left(E_2(z) E_4(z) - E_6(z) - \frac{6i}{\pi z} E_4(z)\right)^2}{\Delta(z)} \\
    &= \frac{(E_2(z) E_4(z) - E_6(z))^{2} - \frac{12i}{\pi z}(E_2(z) E_4(z) - E_6(z)) E_4(z) - \frac{36}{\pi^2 z^2} E_4(z)^{2}}{\Delta(z)} \\
    &= \phi_0(z) - \frac{12 i}{\pi z} \phi_{-2}(z) - \frac{36}{\pi^2 z^2} \phi_{-4}(z).
\end{align}
\end{proof}

% The original expression of $a(x)$ in \cite{Via2017} is given by
% \begin{align}\label{eqn:a-definition-orig}
%     a(x):=&\int\limits_{-1}^i\phi_0\Big(\frac{-1}{z+1}\Big)\,(z+1)^2\,e^{\pi i \|x\|^2 z}\,dz
%     +\int\limits_{1}^i\phi_0\Big(\frac{-1}{z-1}\Big)\,(z-1)^2\,e^{\pi i \|x\|^2 z}\,dz\\
%     -&2\int\limits_{0}^i\phi_0\Big(\frac{-1}{z}\Big)\,z^2\,e^{\pi i \|x\|^2 z}\,dz
%     +2\int\limits_{i}^{i\infty}\phi_0(z)\,e^{\pi i \|x\|^2 z}\,dz.\nonumber
% \end{align}
For our formalization, we choose rectangular contours for the first and the second integral, and write it as follows.
\begin{definition}\label{def:a-definition}\lean{MagicFunction.a.RealIntegrals.a',MagicFunction.a.RadialFunctions.a}\leanok
\uses{def:phi4-phi2-phi0}
Define $a_{\rad} : \R \to \C$ by
\begin{align}\label{eqn:a-rad-definition}
    a_\rad(r) := I_1(r) + I_2(r) + I_3(r) + I_4(r) + I_5(r) + I_6(r)
\end{align}
where
\begin{align}
    I_1(r) &:= \int_{-1}^{-1 + i} \phi_0 \left(\frac{-1}{z+1}\right) (z + 1)^2 e^{\pi i r z} \dd z \label{eqn:a-I1} \\
    I_2(r) &:= \int_{-1 + i}^{i} \phi_0 \left(\frac{-1}{z+1}\right) (z + 1)^2 e^{\pi i r z} \dd z \label{eqn:a-I2} \\
    I_3(r) &:= \int_{1}^{1 + i} \phi_0 \left(\frac{-1}{z-1}\right) (z - 1)^2 e^{\pi i r z} \dd z \label{eqn:a-I3} \\
    I_4(r) &:= \int_{1 + i}^{i} \phi_0 \left(\frac{-1}{z-1}\right) (z - 1)^2 e^{\pi i r z} \dd z \label{eqn:a-I4} \\
    I_5(r) &:= -2 \int_{0}^{i} \phi_0 \left(\frac{-1}{z}\right) z^2 e^{\pi i r z} \dd z \label{eqn:a-I5} \\
    I_6(r) &:= 2 \int_{i}^{i\infty} \phi_0(z) e^{\pi i r z} \dd z \label{eqn:a-I6}
\end{align}
Here all the contours are chosen to be straight line segments.
Now, define $a(x) := a_{\rad}(\|x\|^2)$ for $x \in \R^8$.
\end{definition}
In the original paper, the integrals $I_1$ and $I_2$ are combined, as well as $I_3$ and $I_4$, with the contours running from $\mp 1$ to $i$.
Since the integrands are holomorphic on $\h$ (multiply out \eqref{eqn:phi0-trans-S}), the choice of contours between these endpoints does not matter; we choose rectangular ones because the versions of the Cauchy--Goursat theorem available in Mathlib are stated for rectangles.
In the formalisation, each $I_j$ is realised as an integral over the real parametrising variable $t \in [0,1]$ (or $t \in [1, \infty)$ for $I_6$), so that the \verb|intervalIntegral| API applies.

\subsection{Bounds for the Integrands}

The key estimate is the following bound for quotients by $\Delta$, applicable to all integrands appearing in the definitions of $a$ and $b$.

\begin{lemma}\label{lemma:mod-div-disc-bound}\lean{MagicFunction.PolyFourierCoeffBound.DivDiscBoundOfPolyFourierCoeff}\leanok
Let $f(z)$ be a holomorphic function with a Fourier expansion
\begin{equation}
    f(z) = \sum_{n \ge n_0} c_f(n) e^{\pi i n z}
\end{equation}
with $c_f(n_0) \ne 0$.
Assume that $c_f(n)$ has a polynomial growth, i.e. $|c_f(n)| = O(n^k)$ for some $k \in \N$.
Then there exists a constant $C_f > 0$ such that
\begin{equation}
    \left|\frac{f(z)}{\Delta(z)}\right| \le C_f e^{-\pi (n_0 - 2) \Im z}
\end{equation}
for all $z$ with $\Im z > 1/2$.
\end{lemma}
When $f(z)$ is a (quasi)modular form, we can take $k$ to be the weight of $f$.


\begin{proof}\uses{def:disc-definition}\leanok
By the product formula \eqref{eqn:disc-definition},
\begin{align}
    \left|\frac{f(z)}{\Delta(z)}\right| &= \left|\frac{\sum_{n \ge n_0} c_f(n) e^{\pi i n z}}{e^{2 \pi i z}\prod_{n \ge 1} (1 - e^{2\pi i n z})^{24}}\right| \\
    &= |e^{\pi i (n_0 - 2)z}| \cdot \frac{|\sum_{n \ge n_0} c_f(n) e^{\pi i (n - n_0) z}|}{\prod_{n \ge 1} |1 - e^{2\pi i n z}|^{24}} \\
    &\le e^{-\pi (n_0 - 2) \Im z} \cdot \frac{\sum_{n \ge n_0} |c_f(n)| e^{-\pi (n - n_0) \Im z}}{\prod_{n \ge 1} (1 - e^{- 2\pi n \Im z})^{24}} \\
    &\le e^{-\pi (n_0 - 2) \Im z} \cdot \frac{\sum_{n \ge n_0} |c_f(n)| e^{-\pi (n - n_0) / 2}}{\prod_{n \ge 1} (1 - e^{-\pi n})^{24}} \\
    &= C_f \cdot e^{-\pi (n_0 - 2) \Im z}
\end{align}
where
\begin{equation}
    C_f = \frac{\sum_{n \ge n_0} |c_f(n)| e^{-\pi (n - n_0) / 2}}{\prod_{n \ge 1} (1 - e^{-\pi n})^{24}}.
\end{equation}
The summation in the numerator converges absolutely because of polynomial growth, and the denominator converges to the positive value $e^{\pi} \cdot \Delta(i/2)$ by the product formula. The manipulations of infinite products used here (multiplicativity of the absolute value over convergent products, and monotonicity of convergent products with nonnegative terms) are available in Mathlib.
\end{proof}

To apply Lemma~\ref{lemma:mod-div-disc-bound} to $\phi_0$, $\phi_{-2}$ and $\phi_{-4}$ --- and, later, to the $\psi$-functions --- we need to know that the numerators have Fourier expansions with polynomially growing coefficients, and we need their leading exponents. For a function $f$ with Fourier expansion as in Lemma~\ref{lemma:mod-div-disc-bound}, write $n_0(f)$ for the smallest index with $c_f(n_0(f)) \neq 0$. The following closure property reduces both questions to the corresponding facts about $E_2$, $E_4$, $E_6$ and the $H$-functions.

\begin{lemma}\label{lemma:poly-growth-mul}
Let $f_1, f_2 : \h \to \C$ have absolutely convergent Fourier expansions $f_i(z) = \sum_{n \geq n_0(f_i)} c_i(n)\, e^{\pi i n z}$ with polynomially growing coefficients. Then $f_1 f_2$ has an absolutely convergent Fourier expansion with polynomially growing coefficients, and
\[
  n_0(f_1 f_2) = n_0(f_1) + n_0(f_2), \qquad c_{f_1 f_2}(n_0(f_1 f_2)) = c_1(n_0(f_1))\, c_2(n_0(f_2)) \neq 0.
\]
The analogous statement for $f_1 + f_2$ holds with $n_0(f_1 + f_2) \geq \min(n_0(f_1), n_0(f_2))$ (leading coefficients may cancel).
\end{lemma}
\begin{proof}
By absolute convergence, the product expands as a Cauchy product: the coefficient of $e^{\pi i \ell z}$ in $f_1 f_2$ is $c(\ell) = \sum_{m + n = \ell} c_1(n) c_2(m)$, a sum of at most $\ell + 1$ terms each of size $O(\ell^{k_1 + k_2})$, so $c(\ell) = O(\ell^{k_1 + k_2 + 1})$. The lowest-order term arises exactly once, from $(n, m) = (n_0(f_1), n_0(f_2))$.
\end{proof}

As corollaries, we have the following bounds for $\phi_0$, $\phi_{-2}$, and $\phi_{-4}$.

\begin{corollary}\label{cor:phi0-bound}\uses{thm:ramanujan-formula, lemma:mod-div-disc-bound, lemma:Ek-Fourier, lemma:mod_form_poly_growth, lemma:poly-growth-mul}
There exists a constant $C_0 > 0$ such that
\begin{equation}\label{eqn:phi0-bound}
|\phi_0(z)| \le C_0 e^{-2 \pi \Im z}
\end{equation}
for all $z$ with $\Im z > 1/2$.
\end{corollary}
\begin{proof}
The Fourier coefficients of $E_2$, $E_4$ and $E_6$ grow polynomially (Lemma~\ref{lemma:Ek-Fourier}), so by Lemma~\ref{lemma:poly-growth-mul} the same holds for the numerator $(E_2E_4 - E_6)^2$ of $\phi_0$. By Ramanujan's formula \eqref{eqn:DE4}, $E_2 E_4 - E_6 = 3E_4' = 720 \sum_{n \ge 1} n \sigma_3(n) e^{2 \pi i n z}$, so $n_0(E_2E_4 - E_6) = 2$ in the convention of Lemma~\ref{lemma:mod-div-disc-bound}, and hence
\ifplastex
\begin{equation*}
    (E_2(z) E_4(z) - E_6(z))^{2} = 720^{2} e^{4 \pi i z} + O(e^{5 \pi i z}),
\end{equation*}
\else
\begin{equation}
    (E_2(z) E_4(z) - E_6(z))^{2} = 720^{2} e^{4 \pi i z} + O(e^{5 \pi i z}), \notag
\end{equation}
\fi
i.e.\ $n_0 = 4$. The result follows from Lemma~\ref{lemma:mod-div-disc-bound}.
\end{proof}

\begin{corollary}\label{cor:phi2-bound}\uses{def:phi4-phi2-phi0, lemma:mod-div-disc-bound, lemma:Ek-Fourier, lemma:poly-growth-mul}
There exists a constant $C_{-2} > 0$ such that
\begin{equation}\label{eqn:phi2-bound}
    |\phi_{-2}(z)| \le C_{-2}
\end{equation}
for all $z$ with $\Im z > 1/2$.
\end{corollary}
\begin{proof}
As in Corollary~\ref{cor:phi0-bound}: the numerator of $\phi_{-2}$ is $E_4(E_2E_4 - E_6)$, with $n_0 = 0 + 2 = 2$, so Lemma~\ref{lemma:mod-div-disc-bound} gives the bound with exponent $-\pi(2-2)\Im z = 0$.
\end{proof}

\begin{corollary}\label{cor:phi4-bound}\uses{def:phi4-phi2-phi0, lemma:mod-div-disc-bound, lemma:Ek-Fourier, lemma:poly-growth-mul}
There exists a constant $C_{-4} > 0$ such that
\begin{equation}\label{eqn:phi4-bound}
    |\phi_{-4}(z)| \le C_{-4} e^{2 \pi \Im z}
\end{equation}
for all $z$ with $\Im z > 1/2$.
\end{corollary}
\begin{proof}
The numerator of $\phi_{-4}$ is $E_4^2$, with $n_0 = 0$; apply Lemma~\ref{lemma:mod-div-disc-bound}.
\end{proof}

Note that we can take the constants $C_0$, $C_{-2}$, and $C_{-4}$ as
\begin{align}
    C_0 &= 9 \cdot 240^2 \cdot e^{\pi} \cdot \frac{E_4'(i/2)^{2}}{\Delta(i/2)} \notag \\
    C_{-2} &= 3 \cdot \frac{E_4(i/2) E_4'(i/2)}{\Delta(i/2)} \notag \\
    C_{-4} &= e^{-\pi} \cdot \frac{E_4(i/2)^{2}}{\Delta(i/2)}. \notag
\end{align}

\subsection{Decay of the Integrals $I_j$}

\begin{lemma}\label{lem:bound-I1-I3-I5}\uses{def:a-definition, cor:phi0-bound}\lean{MagicFunction.a.IntegralEstimates.I₁.I₁'_bounding, MagicFunction.a.IntegralEstimates.I₃.I₃'_bounding, MagicFunction.a.IntegralEstimates.I₅.I₅'_bounding}\leanok
    There exists $C > 0$ such that for all $r \geq 0$,
    \begin{equation}
        |I_1(r)|, |I_3(r)|, |I_5(r)| \leq C \int_1^{\infty} e^{-2\pi s} \, e^{-\pi r / s} \, \dd s.
    \end{equation}
\end{lemma}
\begin{proof}\leanok
    We only prove the bound for $I_1(r)$, as the other two are similar.
    By the change of variable $z = -1 + i t$ for $t \in [0,1]$, we have
    $$
    I_1(r) = -i \int_0^1 \phi_0\left(\frac{-1}{i t}\right) t^2 e^{-\pi i r} \cdot e^{-\pi r t} \, \dd t.
    $$
    With the change of variable $s = 1 / t$ (an application, in the formalisation, of Mathlib's one-dimensional change-of-variables theorem \verb|MeasureTheory.integral_image_eq_integral_abs_deriv_smul|), we get
    $$
    I_1(r) = i \int_1^{\infty} \phi_0(i s) \cdot s^{-4} \cdot e^{-\pi i r} \cdot e^{-\pi r / s} \, \dd s
    $$
    and, by the triangle inequality and monotonicity of the integral,
    $$
    |I_1(r)| \leq \int_1^{\infty} |\phi_0(i s)| \cdot s^{-4} \cdot |e^{-\pi i r}| \cdot e^{-\pi r / s} \, \dd s \le \int_1^{\infty} |\phi_0(is)| \cdot e^{-\pi r / s} \, \dd s.
    $$
    Now, Corollary~\ref{cor:phi0-bound} shows
    $$
    |I_1(r)| \leq C_0 \int_1^{\infty} e^{-2\pi s} \, e^{-\pi r / s} \, \dd s.
    $$
\end{proof}

For the integrals $I_2$, $I_4$, and $I_6$, whose contours do not touch the real axis, the decay is exponential.
\begin{lemma}\label{lem:bound-I2-I4-I6}\uses{def:a-definition, cor:phi0-bound}\lean{MagicFunction.a.IntegralEstimates.I₂.I₂'_bounding, MagicFunction.a.IntegralEstimates.I₄.I₄'_bounding, MagicFunction.a.IntegralEstimates.I₆.I₆'_bounding}\leanok
    There exist $C_1, C_2 > 0$ such that for all $r \geq 0$,
    \begin{equation}
        |I_2(r)|, |I_4(r)| \leq C_1 e^{-\pi r}
    \end{equation}
    and
    \begin{equation}
        |I_6(r)| \leq C_2 \frac{e^{-\pi(r + 2)}}{r + 2}
    \end{equation}
\end{lemma}
\begin{proof}
    For $I_2(r)$, parametrize $z$ as $z = t + i$ for $t \in [-1,0]$, and we have
    $$
        I_2(r) = \int_{-1}^0 \phi_0\left(\frac{-1}{t + 1 + i}\right) (t + 1 + i)^2 e^{\pi i r t} e^{-\pi r} \, \dd t.
    $$
    Since the function $z \mapsto \phi_0\left(\frac{-1}{z+1}\right) (z+1)^2$ is holomorphic and bounded on the contour,
    there exists $C > 0$ such that $|\phi_0\left(\frac{-1}{z+1}\right) (z + 1)^2| \leq C$, and the absolute value of the integral can be bounded by
    $$
        |I_2(r)| \le \int_{-1}^{0} C e^{-\pi r} \, \dd t = C e^{-\pi r}.
    $$
    The bound for $I_4(r)$ is similar.
    For $I_6(r)$, parametrize $z$ as $z = i t$ for $t \in [1, \infty)$, and we have
    $$
        I_6(r) = 2 i \int_1^{\infty} \phi_0(i t) e^{-\pi r t} \, \dd t
    $$
    and the absolute value can be bounded as (using Corollary~\ref{cor:phi0-bound})
    $$
        |I_6(r)| \leq 2 \int_1^{\infty} |\phi_0(i t)| e^{-\pi r t} \, \dd t \leq 2C_0 \int_1^{\infty} e^{-2\pi t} e^{-\pi r t} \, \dd t = \frac{2C_0}{\pi} \frac{e^{-\pi (r + 2)}}{r + 2}.
    $$
\end{proof}

The bound appearing in Lemma~\ref{lem:bound-I1-I3-I5} decays faster than any polynomial as $r \to \infty$; this is a consequence of standard facts about the $\Gamma$ function.

\begin{lemma}\label{lem:integral-bound}
    For all $n \in \N$, there exists a constant $C'$ such that for all $r \geq 0$,
    \begin{align}
        r^n \cdot \int_{1}^{\infty} e^{-2\pi s} \, e^{-\pi r /s} \, \dd s \leq C'
    \end{align}
\end{lemma}
\begin{proof}
    Fix $n \in \N$. We know there exists a constant $C$ such that for all $x \geq 0$, $\abs{x}^n \cdot \abs{e^{-\pi x}} \leq C$. In particular, for all $r \geq 0$ and $s \geq 1$, $r^n \cdot e^{-\pi r/s} \leq C s^n$. Hence, for all $r \geq 0$, we can write
    \begin{align}
        r^n \cdot \int_{1}^{\infty} e^{-2\pi s} \, e^{-\pi r /s} \, \dd{s}
        = \int_{1}^{\infty} e^{-2\pi s} \, \left(\abs{r}^n \cdot e^{-\pi r /s}\right) \, \dd{s}
        \leq C \int_{1}^{\infty} e^{-2\pi s} \, s^n \, \dd{s}
    \end{align}
    The $\Gamma$ function is given by
    \begin{align}
        \Gamma(x) = \int_{0}^{\infty} e^{-u} \, u^{x-1} \, \dd{u}
    \end{align}
    for all $x > 0$ (in Mathlib, \verb|Real.Gamma_eq_integral|). Hence, writing $u = 2\pi s$ and using $\Gamma(n+1) = n!$ (\verb|Real.Gamma_nat_eq_factorial|), we get
    \begin{align}
        C \int_{1}^{\infty} e^{-2\pi s} \, s^n \, \dd{s}
        \leq C \int_{0}^{\infty} e^{-2\pi s} \, s^n \, \dd{s}
        = C \int_{0}^{\infty} \frac{1}{(2\pi)^{n+1}} e^{-u} \, u^n \, \dd{u}
        = \frac{C}{(2\pi)^n} \Gamma(n+1)
        = \frac{C \cdot n!}{(2\pi)^n}
    \end{align}
    Defining $C' := \frac{C \cdot n!}{(2\pi)^n}$ finishes the proof.
\end{proof}

\subsection{Smoothness of the Integrals $I_j$, and Schwartzness of $a$}

\begin{lemma}\label{lem:Ij-smooth}\uses{def:a-definition, cor:phi0-bound}
    The functions $I_1, \ldots, I_5$ are smooth on $\R$, and $I_6$ is smooth on $(-2, \infty)$. Moreover, for every $j$ and $n$, the $n$-th derivative $I_j^{(n)}$ is given by inserting the factor $(\pi i z)^n$ into the corresponding integrand, and consequently, for all $k, n \in \N$ there exists $C$ such that $r^k\, |I_j^{(n)}(r)| \leq C$ for all $r \geq 0$.
\end{lemma}
\begin{proof}\uses{lem:bound-I1-I3-I5, lem:bound-I2-I4-I6, lem:integral-bound}
    Each $I_j$ is an integral of $g_j(z)\, e^{\pi i r z}$ over a fixed contour, with $g_j$ holomorphic; formally, over the real parametrising variable. For $j \leq 5$ the contour is compact and the integrand, together with each of its $r$-derivatives $g_j(z) (\pi i z)^n e^{\pi i r z}$, is dominated locally uniformly in $r$ by an integrable function; the Leibniz integral rule (in Mathlib, \verb|hasDerivAt_integral_of_dominated_loc_of_deriv_le|) therefore applies repeatedly and yields smoothness on $\R$. For $I_6$, Corollary~\ref{cor:phi0-bound} bounds the integrand by a multiple of $e^{-\pi(r+2)t}$ on the contour $z = it$, $t \geq 1$; this is integrable precisely when $r > -2$, and dominates all $r$-derivatives locally uniformly on $(-2, \infty)$, giving smoothness there. (For $r \leq -2$ the integral defining $I_6$ diverges; this is what makes the cutoff of Corollary~\ref{cor:schwartz-of-nonneg-decay} necessary.)

    For the derivative bounds on $[0, \infty)$: the extra factor $(\pi i z)^n$ is bounded on the compact contours, and on the unbounded contours (after the substitution $s = 1/t$ in the case of $I_1, I_3, I_5$) contributes a factor bounded by $(\pi\sqrt{2})^n$ resp.\ a polynomial in $t$, which is absorbed by the exponential decay of the integrand. The bounds of Lemmas~\ref{lem:bound-I1-I3-I5} and~\ref{lem:bound-I2-I4-I6} therefore hold, with adjusted constants, for all derivatives, and Lemma~\ref{lem:integral-bound} converts them into the stated rapid decay.
\end{proof}

The functions $I_j$ are \emph{not} Schwartz functions on $\R$: the bounds above degrade into exponential \emph{growth} as $r \to -\infty$, and $I_6$ is not even defined by a convergent integral for $r \leq -2$. This is harmless: $a(x) = a_\rad(\|x\|^2)$ only sees $a_\rad$ on $[0, \infty)$, and Lemma~\ref{lem:Ij-smooth} shows that $a_\rad$ satisfies the hypotheses of Corollary~\ref{cor:schwartz-of-nonneg-decay} with $\theta = -2$.

\begin{proposition}\label{prop:a-schwartz}\lean{MagicFunction.FourierEigenfunctions.a}\leanok
\uses{def:a-definition, cor:phi0-bound}
$a(x)$ is a Schwartz function.
\end{proposition}
\begin{proof}
\uses{lem:bound-I1-I3-I5, lem:bound-I2-I4-I6, lem:integral-bound, lem:Ij-smooth, cor:schwartz-of-nonneg-decay}
By Lemma~\ref{lem:Ij-smooth}, $a_\rad = I_1 + \cdots + I_6$ is smooth on $(-2, \infty)$ and all its derivatives decay rapidly on $[0, \infty)$. Corollary~\ref{cor:schwartz-of-nonneg-decay} (with $\theta = -2$) shows that $x \mapsto a_\rad(\|x\|^2) = a(x)$ is a Schwartz function on $\R^8$.

In the formalisation, this is implemented by realising each product $\Theta_{-2} \cdot I_j$ as a term of type $\mathcal{S}(\R, \C)$, transporting it to $\mathcal{S}(\R^8, \C)$ via Lemma~\ref{lemma:schwartz-comp-normsq}, and summing; since $\Theta_{-2} = 1$ on $[0, \infty)$, the resulting Schwartz function is equal to the function $a$ of Definition~\ref{def:a-definition}.
\end{proof}

\subsection{The Eigenfunction Property of $a$}

We now prove that $a$ satisfies condition \eqref{eqn:a-fourier}. The Fourier transform acts on the summands of $a$ by permuting them: writing (by a slight abuse of notation) $I_j$ also for the radial function $x \mapsto I_j(\|x\|^2)$ on $\R^8$, we shall show
\begin{align}
    \mathcal{F}(I_1 + I_2) &= I_3 + I_4, \label{eqn:a-perm-12}\\
    \mathcal{F}(I_5) &= I_6. \label{eqn:a-perm-56}
\end{align}

\begin{lemma}\label{lemma:perm-I12}\uses{def:a-definition, lemma:Gaussian-Fourier, lemma:phi0-transform, prop:a-schwartz}\lean{MagicFunction.a.Fourier.perm_I₁_I₂}\leanok
\eqref{eqn:a-perm-12} holds: $\mathcal{F}(I_1 + I_2) = I_3 + I_4$.
\end{lemma}
\begin{proof}
Fix $\xi \in \R^8$ and write $F_j := \mathcal{F}(I_j(\|\cdot\|^2))$. The double integral defining $F_1(\xi)$ converges absolutely --- the estimates of Lemmas~\ref{lem:bound-I1-I3-I5} and~\ref{lem:bound-I2-I4-I6} bound the inner integral uniformly by a Schwartz-type bound in $x$ --- so by Fubini's theorem we may exchange the Fourier integral with the contour integral. Lemma~\ref{lemma:Gaussian-Fourier} then gives
\begin{align}
    F_1(\xi) &= \int_{-1}^{-1+i} \phi_0\left(\frac{-1}{z+1}\right) (z+1)^2\, z^{-4}\, e^{\pi i \|\xi\|^2 (\frac{-1}{z})}\, \dd z,
\end{align}
and similarly for $F_2$ with contour from $-1 + i$ to $i$. We substitute $w = -1/z$, which maps the two segments onto circular arcs from $1$ to $\frac{1+i}{2}$ and from $\frac{1+i}{2}$ to $i$ respectively, and compute, using the $1$-periodicity \eqref{eqn:phi0-trans-T} of $\phi_0$,
\[
    \phi_0\left(-1-\frac{1}{w-1}\right)\left(\frac{-1}{w}+1\right)^2 w^2 = \phi_0\left(\frac{-1}{w-1}\right)(w-1)^2,
\]
so that $F_1(\xi) + F_2(\xi)$ is the integral of $\phi_0\left(\frac{-1}{w-1}\right)(w-1)^2 e^{\pi i \|\xi\|^2 w}$ along the concatenated arc from $1$ to $i$. This integrand is holomorphic on $\h$ (multiply out \eqref{eqn:phi0-trans-S}) and continuous up to the boundary, so by the Cauchy--Goursat theorem (in the formalisation: the Mathlib version for rectangles, together with a version for the regions between a quarter-circle and the two sides of the enclosing square) the arc may be replaced by the rectangular path $1 \to 1 + i \to i$. The two resulting integrals are precisely $I_3(\|\xi\|^2)$ and $I_4(\|\xi\|^2)$.
\end{proof}

\begin{lemma}\label{lemma:perm-I56}\uses{def:a-definition, lemma:Gaussian-Fourier, lemma:phi0-transform, prop:a-schwartz}\lean{MagicFunction.a.Fourier.perm_I₅}\leanok
\eqref{eqn:a-perm-56} holds: $\mathcal{F}(I_5) = I_6$.
\end{lemma}
\begin{proof}
As in Lemma~\ref{lemma:perm-I12}, exchange the integrals and apply Lemma~\ref{lemma:Gaussian-Fourier}, then substitute $w = -1/z$. Here no contour deformation is needed: the substitution maps the segment from $0$ to $i$ onto the ray from $i\infty$ to $i$, reversing orientation, and
\[
    -2\,\phi_0\left(\frac{-1}{z}\right) z^2\, z^{-4}\, e^{\pi i \|\xi\|^2(\frac{-1}{z})}\,\dd z \;\longmapsto\; 2\,\phi_0(w)\, e^{\pi i \|\xi\|^2 w}\,\dd w,
\]
which is the integrand of $I_6$.
\end{proof}

\begin{lemma}\label{lemma:perm-I-reverse}\uses{lemma:perm-I12, lemma:perm-I56, lemma:radial-fourier-inversion}\lean{MagicFunction.a.Fourier.perm_₃_I₄, MagicFunction.a.Fourier.perm_I₆}\leanok
$\mathcal{F}(I_3 + I_4) = I_1 + I_2$ and $\mathcal{F}(I_6) = I_5$.
\end{lemma}
\begin{proof}\leanok
The functions $I_1 + I_2$ and $I_5$ (as functions on $\R^8$) are radial Schwartz functions, so $\mathcal{F}^2$ fixes them (Lemma~\ref{lemma:radial-fourier-inversion}). Applying $\mathcal{F}$ to \eqref{eqn:a-perm-12} and \eqref{eqn:a-perm-56} gives the claim.
\end{proof}

\begin{proposition}\label{prop:a-fourier}
\uses{def:a-definition, prop:a-schwartz, lemma:perm-I12, lemma:perm-I56, lemma:perm-I-reverse}\lean{MagicFunction.a.Fourier.eig_a}\leanok
$a(x)$ satisfies \eqref{eqn:a-fourier}.
\end{proposition}
\begin{proof}\leanok
By linearity of $\mathcal{F}$ on the Schwartz space,
\[
    \mathcal{F}(a) = \mathcal{F}(I_1 + I_2) + \mathcal{F}(I_3 + I_4) + \mathcal{F}(I_5) + \mathcal{F}(I_6)
    = (I_3 + I_4) + (I_1 + I_2) + I_6 + I_5 = a
\]
by Lemmas~\ref{lemma:perm-I12}, \ref{lemma:perm-I56} and \ref{lemma:perm-I-reverse}.
\end{proof}

\subsection{Double Zeroes of $a$}

Next, we check that $a$ has double zeroes at all $\Lambda_8$-lattice points of length greater then $\sqrt{2}$.
The contour deformations in this argument involve unbounded contours; the following consequence of the Cauchy--Goursat theorem, formalised as part of this project, is what we need.

\begin{lemma}[Cauchy--Goursat for unbounded rectangular contours]\label{lem:CG-unbounded}\lean{Complex.tendsto_integral_boundary_open_rect_one_side_atTop_nhds_sum_other_two_sides, Complex.integral_boundary_open_rect_eq_zero_of_differentiable_on_off_countable}\leanok
Let $f : \C \to \C$ be continuous on the closed strip $\{x_1 \leq \Re z \leq x_2,\ \Im z \geq y\}$, holomorphic on its interior, with $f(z) \to 0$ as $\Im z \to \infty$ (uniformly in $\Re z$), and integrable on the two vertical rays. Then
\[
    \int_{x_1 + iy}^{x_1 + i\infty} f(z)\, \dd z
    = \int_{x_1 + iy}^{x_2 + iy} f(z)\, \dd z + \int_{x_2 + iy}^{x_2 + i\infty} f(z)\, \dd z.
\]
\end{lemma}
\begin{proof}\leanok
Apply the Cauchy--Goursat theorem for rectangles (available in Mathlib) to the truncated rectangles with upper edge at height $T$, and let $T \to \infty$: the integral over the upper edge tends to $0$ by the decay hypothesis, and the integrals over the truncated vertical sides tend to the integrals over the rays by integrability.
\end{proof}

Using \eqref{eqn:phi0-bound}, \eqref{eqn:phi2-bound}, and \eqref{eqn:phi4-bound}, we can control the behavior of $\phi_0$ near $0$ and $i\infty$.

\begin{corollary}\label{cor:phi0-near-0-infty}\uses{cor:phi0-bound, cor:phi2-bound, cor:phi4-bound, lemma:phi0-transform}
We have
\begin{align}
    \phi_0\left(\frac{i}{t}\right) &= O(e^{-2 \pi / t}) \quad \text{as } t \to 0 \label{eqn:phi0-near-0} \\
    \phi_0\left(\frac{i}{t}\right) &= O(t^{-2}e^{2 \pi t}) \quad \text{as } t \to \infty. \label{eqn:phi0-near-infty} \\
\end{align}
\end{corollary}
\begin{proof}
The first estimate follows from \eqref{eqn:phi0-bound} with $z = i/t$.
For the second estimate, by \eqref{eqn:phi0-trans-S}, \eqref{eqn:phi2-bound}, and \eqref{eqn:phi4-bound}, we have
\begin{equation}
    \left|\phi_0\left(\frac{i}{t}\right)\right| \le |\phi_0(it)| + \frac{12}{\pi t} |\phi_{-2}(it)| + \frac{36}{\pi^2 t^2} |\phi_{-4}(it)|
    \le C_0 e^{-2 \pi t} + \frac{12}{\pi t} \cdot C_{-2} + \frac{36}{\pi^2 t^2} \cdot C_{-4} e^{2 \pi t} = O(t^{-2}e^{2 \pi t}).
\end{equation}
\end{proof}

In particular, for $r > 2$ the integrand $\phi_0(-1/z)\, z^2\, e^{\pi i r z}$ is bounded by a multiple of $e^{\pi (2 - r)\Im z}$ for $\Im z \geq 2$, so it tends to $0$ as $\Im(z) \to \infty$ and is integrable on vertical rays; the same applies to its translates by $\pm 1$. This provides the hypotheses of Lemma~\ref{lem:CG-unbounded} in the following proof. (In the radial variable, these are inputs of the form $r = \|x\|^2 > 2$, i.e.\ $\|x\| > \sqrt 2$.)

\begin{proposition}
\label{prop:a-double-zeros}
\uses{cor:phi0-near-0-infty, def:a-definition, cor:disc-nonvanishing, lem:CG-unbounded}
For $r>\sqrt{2}$ we can express $a(r)$ in the following form
\begin{equation}\label{eqn: a double zeroes}
    a(r)=-4\sin(\pi r^2/2)^2\,\int\limits_{0}^{i\infty}\phi_0\Big(\frac{-1}{z}\Big)\,z^2\,e^{\pi i r^2 \,z}\,dz.
\end{equation}
\end{proposition}
\begin{proof}
We denote the right hand side of \eqref{eqn: a double zeroes} by $d(r)$.
Convergence of the integral for $r > \sqrt{2}$ follows from Corollary~\ref{cor:phi0-near-0-infty}.
Writing $-4\sin(\pi r^2/2)^2 = e^{\pi i r^2} - 2 + e^{-\pi i r^2}$ and absorbing the exponentials into shifts $z \mapsto z \pm 1$ of the integration variable, we can write
\begin{align}
    d(r)=&\int\limits_{-1}^{i\infty-1}\phi_0\Big(\frac{-1}{z+1}\Big)\,(z+1)^2\,e^{\pi i r^2 \,z}\,dz-
    2\int\limits_{0}^{i\infty}\phi_0\Big(\frac{-1}{z}\Big)\,z^2\,e^{\pi i r^2 \,z}\,dz\notag\\
    +&\int\limits_{1}^{i\infty+1}\phi_0\Big(\frac{-1}{z-1}\Big)\,(z-1)^2\,e^{\pi i r^2 \,z}\,dz.\notag
\end{align}
From \eqref{eqn:phi0-trans-S} we deduce that if $r>\sqrt{2}$ then
$\phi_0\Big(\frac{-1}{z}\Big)\,z^2\,e^{\pi i r^2 \,z}\to 0$ as $\Im(z)\to\infty$. Therefore, we can deform the paths of integration by Lemma~\ref{lem:CG-unbounded}
and rewrite
\begin{align}
    d(r)=&\int\limits_{-1}^{i}\phi_0\Big(\frac{-1}{z+1}\Big)\,(z+1)^2\,e^{\pi i r^2 \,z}\,dz
    +\int\limits_{i}^{i\infty}\phi_0\Big(\frac{-1}{z+1}\Big)\,(z+1)^2\,e^{\pi i r^2 \,z}\,dz\notag\\
    -2&\int\limits_{0}^{i}\phi_0\Big(\frac{-1}{z}\Big)\,z^2\,e^{\pi i r^2 \,z}\,dz
    -2\int\limits_{i}^{i\infty}\phi_0\Big(\frac{-1}{z}\Big)\,z^2\,e^{\pi i r^2 \,z}\,dz\notag\\
    +&\int\limits_{1}^{i}\phi_0\Big(\frac{-1}{z-1}\Big)\,(z-1)^2\,e^{\pi i r^2 \,z}\,dz
    +\int\limits_{i}^{i\infty}\phi_0\Big(\frac{-1}{z-1}\Big)\,(z-1)^2\,e^{\pi i r^2 \,z}\,dz.\notag
\end{align}
Now from \eqref{eqn:phi0-trans-S} we find
\begin{align}&\phi_0\Big(\frac{-1}{z+1}\Big)\,(z+1)^2-2\phi_0\Big(\frac{-1}{z}\Big)\,z^2+
\phi_0\Big(\frac{-1}{z-1}\Big)\,(z-1)^2=\notag\\
&\phi_0(z+1)\,(z+1)^2-2\phi_0(z)\,z^2+\phi_0(z-1)\,(z-1)^2\notag\\
&-\frac{12i}{\pi}\,\Big(\phi_{-2}(z+1)\,(z+1)-2\phi_{-2}(z)\,z+\phi_{-2}(z-1)\,(z-1)\Big)\notag\\
&-\frac{36}{\pi^2}\Big(\phi_{-4}(z+1)-2\phi_{-4}(z)+\phi_{-4}(z-1)\Big)=\notag\\
&2\phi_0(z),
\end{align}
using the $1$-periodicity of $\phi_0$, $\phi_{-2}$, $\phi_{-4}$. Thus, we obtain
\begin{align}
    d(r)=&\int\limits_{-1}^{i}\phi_0\Big(\frac{-1}{z+1}\Big)\,(z+1)^2\,e^{\pi i r^2 \,z}\,dz
    -2\int\limits_{0}^{i}\phi_0\Big(\frac{-1}{z}\Big)\,z^2\,e^{\pi i r^2 \,z}\,dz\notag\\
    +&\int\limits_{1}^{i}\phi_0\Big(\frac{-1}{z-1}\Big)\,(z-1)^2\,e^{\pi i r^2 \,z}\,dz
    +2\int\limits_{i}^{i\infty}\phi_0(z)\,e^{\pi i r^2 \,z}\,dz=a(r).\notag
\end{align}
This finishes the proof.
\end{proof}

The factor $\sin(\pi r^2/2)^2$ vanishes to second order at $r = \sqrt{2n}$ for every integer $n$, so Proposition~\ref{prop:a-double-zeros} exhibits the required double zeroes of $a$ at all $\Lambda_8$-vectors of length greater than $\sqrt 2$.

\subsection{Special Values of $a$}

Finally, we find another convenient integral representation for $a$ and compute the value of $a$ at $0$.

\begin{proposition}\label{prop:a-another-integral}\uses{prop:a-double-zeros, def:phi4-phi2-phi0, lemma:phi0-transform, def:a-definition}
For $r\geq0$ we have
\begin{align}\label{eqn:a-another-integral}a(r)=&4i\,\sin(\pi r^2/2)^2\,\Bigg(\frac{36}{\pi^3\,(r^2-2)}-\frac{8640}{\pi^3\,r^4}+\frac{18144}{\pi^3\,r^2}\\ +&\int\limits_0^\infty\,\left(t^2\,\phi_0\Big(\frac{i}{t}\Big)-\frac{36}{\pi^2}\,e^{2\pi t}+\frac{8640}{\pi}\,t-\frac{18144}{\pi^2}\right)\,e^{-\pi r^2 t}\,dt \Bigg) .\notag\end{align}
The integral converges absolutely for all $r\in\R_{\geq 0}$.
\end{proposition}
\begin{proof}
Suppose that $r>\sqrt{2}$. Then by Proposition~\ref{prop:a-double-zeros}
$$a(r)=4i\,\sin(\pi r^2/2)^2\,\int\limits_{0}^{\infty}\phi_0(i/t)\,t^2\,e^{-\pi r^2 t}\,dt. $$
From \eqref{eqn:phi0-trans-S} we obtain
\begin{equation}\label{eqn: phi asymptotic}
\phi_0(i/t)\,t^2=\frac{36}{\pi^2}\,e^{2 \pi t}-\frac{8640}{\pi}\,t+\frac{18144}{\pi^2}+O(t^2\,e^{-2\pi t})\quad\mbox{as}\;t\to\infty.
\end{equation}
For $r>\sqrt{2}$ we have
\begin{equation}
\int\limits_0^\infty \left(\frac{36}{\pi^2}\,e^{2 \pi t}+\frac{8640}{\pi}\,t+\frac{18144}{\pi^2}\right)\,e^{-\pi r^2 t}\,dt
=\frac{36}{\pi^3\,(r^2-2)}-\frac{8640}{\pi^3\,r^4}+\frac{18144}{\pi^3\,r^2}.\end{equation}
Therefore, the identity \eqref{eqn:a-another-integral} holds for $r>\sqrt{2}$.

On the other hand, from the definition~\eqref{eqn:a-rad-definition} we see that $a(r)$ is analytic in some neighborhood of $[0,\infty)$. The asymptotic expansion~\eqref{eqn: phi asymptotic} implies that the right hand side of \eqref{eqn:a-another-integral} is also analytic in some neighborhood of $[0,\infty)$. Hence, by the identity principle for analytic functions (available in Mathlib as \verb|AnalyticOnNhd.eqOn_of_preconnected_of_eventuallyEq|), the identity \eqref{eqn:a-another-integral} holds on the whole interval $[0,\infty)$. This finishes the proof of the proposition.
\end{proof}
From the identity~\eqref{eqn:a-another-integral} we see that the values $a(r)$ are in $i\R$ for all $r\in\R_{\geq0}$.
% In particular, we have
\begin{proposition}\label{prop:a0}\uses{prop:a-another-integral}\lean{MagicFunction.a.SpecialValues.a_zero}\leanok
We have $a(0) = \frac{-8640\, i}{\pi}$.
\end{proposition}
\begin{proof}
As $r \to 0$, the factor $\sin(\pi r^2/2)^2$ vanishes to order $r^4$; the only term inside the parentheses of \eqref{eqn:a-another-integral} that blows up at that order is $-\frac{8640}{\pi^3 r^4}$. Hence
\[
  a(0) = \lim_{r \to 0} 4i\left(\frac{\pi r^2}{2}\right)^2 \cdot \left(-\frac{8640}{\pi^3 r^4}\right) = \frac{-8640\, i}{\pi}.
\]
\end{proof}

\subsection{The Integrands for $b$}

Now we construct function $b$. To this end we consider the function
\begin{definition}\label{def: h}\uses{def:H2-H3-H4}\lean{h}\leanok
% \begin{equation}\label{eqn: h define}
%     h(z)\,:=\,128 \frac{\theta_{00}^4(z)+\theta_{01}^4(z)}{\theta_{10}^8(z)}.
% \end{equation}
\begin{equation}\label{eqn: h define}
    h(z) := 128 \frac{H_3(z) + H_4(z)}{H_2(z)^2}.
\end{equation}
\end{definition}
It is easy to see that $h\in M^!_{-2}(\Gamma_0(2))$. Indeed, first we check that $h|_{-2}\gamma=h$
for all $\gamma\in\Gamma_0(2)$. Since the group $\Gamma_0(2)$ is generated by elements
$\left(\begin{smallmatrix}1&0\\2&1\end{smallmatrix}\right)$ and $\left(\begin{smallmatrix}1&1\\0&1\end{smallmatrix}\right)$
it suffices to check that $h$ is invariant under their action. This follows immediately
from \eqref{eqn:H2-transform-S}--\eqref{eqn:H4-transform-S} and \eqref{eqn: h define}. Next we analyze the poles of $h$.
It is known \cite[Chapter~I Lemma~4.1]{Mumford} that $\theta_{10}$ has no zeros in the upper-half plane and hence $h$ has poles only at the cusps.
At the cusp $i\infty$ this modular form has the Fourier expansion
\ifplastex
\begin{equation*}
h(z)\,=\,q^{-1} + 16 - 132 q + 640 q^2 - 2550 q^3+O(q^4).
\end{equation*}
\else
\begin{equation}
h(z)\,=\,q^{-1} + 16 - 132 q + 640 q^2 - 2550 q^3+O(q^4).\notag
\end{equation}
\fi
Let $I=\left(\begin{smallmatrix}1&0\\0&1\end{smallmatrix}\right)$,
$T=\left(\begin{smallmatrix}1&1\\0&1\end{smallmatrix}\right)$, and
$S=\left(\begin{smallmatrix}0&-1\\1&0\end{smallmatrix}\right)$ be elements of $\Gamma_1$.
\begin{definition}\label{def:psiI-psiT-psiS}\uses{def: h}\lean{ψI, ψT, ψS}\leanok
We define the following three functions
\begin{align}
    \psi_I\,:=\,&h-h|_{-2}ST \label{eqn:psiI-define}\\
    \psi_T\,:=\,&\psi_I|_{-2}T \label{eqn:psiT-define}\\
    \psi_S\,:=\,&\psi_I|_{-2}S. \label{eqn:psiS-define}
\end{align}
\end{definition}

The weight $2$ transformation laws of the $H$-functions (Lemma~\ref{lemma:theta-transform-S-T}) yield explicit expressions for the $\psi$-functions.

\begin{lemma}\label{lemma:psi-explicit}\uses{def:psiI-psiT-psiS, lemma:theta-transform-S-T}\lean{ψI_eq, ψT_eq, ψS_eq'}\leanok
\begin{align}
    \psi_I &= 128 \left(\frac{H_3 + H_4}{H_2^2} + \frac{H_4 - H_2}{H_3^2}\right) \label{eqn:psiI-explicit}\\
    \psi_T &= 128 \left(\frac{H_3 + H_4}{H_2^2} + \frac{H_2 + H_3}{H_4^2}\right) \label{eqn:psiT-explicit}\\
    \psi_S &= 128 \left(\frac{H_4 - H_2}{H_3^2} - \frac{H_2 + H_3}{H_4^2}\right) \label{eqn:psiS-explicit}
\end{align}
\end{lemma}
\begin{proof}\leanok
By Lemma \ref{lemma:theta-transform-S-T} and the cocycle property of the slash action,
\begin{align}
    H_2|_{2}ST = -H_3, \qquad
    H_3|_{2}ST = -H_4, \qquad
    H_4|_{2}ST = H_2,
\end{align}
so, since the automorphy factors cancel in the quotient defining $h$,
\[
    h|_{-2}ST = 128\,\frac{(H_3|_2 ST) + (H_4|_2 ST)}{((H_2|_2 ST))^2}
    = 128\, \frac{-H_4 + H_2}{H_3^2},
\]
which gives \eqref{eqn:psiI-explicit}. Applying $|_{-2}T$ and $|_{-2}S$ to $\psi_I$, using Lemma~\ref{lemma:theta-transform-S-T} again, gives \eqref{eqn:psiT-explicit} and \eqref{eqn:psiS-explicit}.
\end{proof}

The same computations give the following relations, which we use repeatedly; all follow from Definition~\ref{def:psiI-psiT-psiS}, Lemma~\ref{lemma:theta-transform-S-T} and the cocycle property of the slash action.

\begin{lemma}\label{lemma:psi-slash-relations}\uses{def:psiI-psiT-psiS, lemma:psi-explicit}\lean{ψT_slash_T, ψS_slash_S, ψT_slash_S, ψS_slash_T, ψS_slash_ST, ψS_slash_TSTS}\leanok
We have
\begin{align}
    \psi_T|_{-2}T &= \psi_I, & \text{i.e.}\quad \psi_I(z) &= \psi_T(z+1), \label{eqn:psi-rel-TT}\\
    \psi_S|_{-2}S &= \psi_I, & \psi_I(z) &= z^2\,\psi_S(-1/z), \label{eqn:psi-rel-SS}\\
    \psi_T|_{-2}S &= -\psi_T, & -\psi_T(z) &= z^2\,\psi_T(-1/z), \label{eqn:psi-rel-TS}\\
    \psi_S|_{-2}T &= -\psi_S, & -\psi_S(z) &= \psi_S(z+1), \label{eqn:psi-rel-ST}\\
    \psi_S|_{-2}ST &= \psi_T, & \psi_T(z) &= (z+1)^2\,\psi_S\!\left(\frac{-1}{z+1}\right), \label{eqn:psi-rel-SST}\\
    \psi_S|_{-2}TSTS &= \psi_T, & \psi_T(z) &= (z-1)^2\,\psi_S\!\left(\frac{-1}{z-1}\right). \label{eqn:psi-rel-STSTS}
\end{align}
\end{lemma}

For the growth estimates it is convenient to express the $\psi$-functions, like the $\phi$-functions, as quotients by $\Delta$.

\begin{lemma}\label{lemma:psi-new}\uses{lemma:psi-explicit, lemma:jacobi-identity, lemma:lv1-lv2-identities}
    $\psi_I(z), \psi_S(z), \psi_T(z)$ can be written as
    \begin{align}
        \psi_I(z) &= \frac{H_4^3 (5 H_2^2 + 5 H_2 H_4 + 2 H_4^2)}{2\Delta}, \label{eqn:psiI-new} \\
        \psi_S(z) &= -\frac{H_2^3 (2 H_2^2 + 5 H_2 H_4 + 5 H_4^2)}{2 \Delta}, \label{eqn:psiS-new} \\
        \psi_T(z) &= \frac{H_3^3 (5 H_2^2 - 5 H_2 H_3 + 2 H_3^2)}{2 \Delta}.\label{eqn:psiT-new}
    \end{align}
    In particular, the $\psi$-functions are holomorphic on $\h$.
\end{lemma}
\begin{proof}
Starting from \eqref{eqn:psiI-explicit} and using the Jacobi identity $H_3 = H_2 + H_4$ \eqref{eqn:jacobi-identity} together with $\Delta = \frac{1}{256}(H_2 H_3 H_4)^2$ \eqref{eqn:disctheta}, we can rewrite $\psi_I(z)$ as
\begin{align}
    \psi_I &= 128\, \frac{H_3^2 (H_3 + H_4) - H_2^2 (H_2 - H_4)}{H_2^2 H_3^2} \\
    &= 128\, \frac{(H_2 + H_4)^2 (H_2 + 2 H_4) - H_2^2 (H_2 - H_4)}{H_2^2 H_3^2} \\
    &= 128\, \frac{H_4(5 H_2^2 + 5 H_2 H_4 + 2 H_4^2)}{ H_2^2 H_3^2}
    = \frac{H_4^3 (5 H_2^2 + 5 H_2 H_4 + 2 H_4^2)}{2\Delta}.
\end{align}
Applying $|_{-2}S$ and $|_{-2}T$ to this expression proves \eqref{eqn:psiS-new} and \eqref{eqn:psiT-new}.
The final statement follows since the $H$-functions are holomorphic and $\Delta$ is non-vanishing on $\h$ (Corollary~\ref{cor:disc-nonvanishing}).
\end{proof}


\begin{lemma}\label{lemma:psiI-psiT-psiS-fourier}\uses{lemma:psi-new, prop:H2-fourier, prop:H3-fourier, prop:H4-fourier}
The Fourier expansions of these functions are
\begin{align}
    \psi_I(z)\,=\,&q^{-1} + 144 + O(q^{1/2}) \label{eqn: psi fourier I}\\
    \psi_T(z)\,=\,&q^{-1} + 144 + O(q^{1/2}) \label{eqn: psi fourier T}
\end{align}
\end{lemma}
\begin{proof}
Direct computation from Lemma~\ref{lemma:psi-new} and the $q$-expansions of the $H$-functions (Propositions~\ref{prop:H2-fourier}--\ref{prop:H4-fourier}).
\end{proof}

Now, we are ready to define function $b$.
Again, the definition here is slightly different from that in \cite{Via2017}, where all the contours are chosen to be the straight lines.


\begin{definition}\label{def:b-definition}\uses{def:psiI-psiT-psiS}\lean{MagicFunction.b.RealIntegrals.b',MagicFunction.b.RadialFunctions.b}\leanok
Define $b_\rad : \R \to \C$ by
\begin{equation}
    \label{eqn:b-definition}
    b_\rad(r) := J_1(r) + J_2(r) + J_3(r) + J_4(r) + J_5(r) + J_6(r)
\end{equation}
where for $r \in \R$,
\begin{align}
    J_1(r) &:= \int_{-1}^{-1 + i} \psi_T(z) e^{\pi i r z} \, \dd z, \label{eqn:J1} \\
    J_2(r) &:= \int_{-1 + i}^{i} \psi_T(z) e^{\pi i r z} \, \dd z, \label{eqn:J2} \\
    J_3(r) &:= \int_{1}^{1 + i} \psi_T(z) e^{\pi i r z} \, \dd z, \label{eqn:J3} \\
    J_4(r) &:= \int_{1 + i}^{i} \psi_T(z) e^{\pi i r z} \, \dd z, \label{eqn:J4} \\
    J_5(r) &:= -2 \int_{0}^{i} \psi_I(z) e^{\pi i r z} \, \dd z, \label{eqn:J5} \\
    J_6(r) &:= -2 \int_{i}^{i \infty} \psi_S(z) e^{\pi i r z} \, \dd z. \label{eqn:J6}
\end{align}
Here all the contours are straight line segments.
Then we define $b : \R^8 \to \C$ by $b(x) := b_\rad(\|x\|^2)$.
\end{definition}

By Lemma~\ref{lemma:psi-slash-relations}, the integrands of $J_1, \ldots, J_5$ can be rewritten in a form exactly parallel to those of $I_1, \ldots, I_5$, with $\psi_S$ playing the role of $\phi_0$:
\begin{align}
    J_1(r) &= \int_{-1}^{-1 + i} \psi_S\!\left(\frac{-1}{z + 1}\right) (z + 1)^2\, e^{\pi i r z} \, \dd z,
    & J_3(r) &= \int_{1}^{1 + i} \psi_S\!\left(\frac{-1}{z - 1}\right) (z - 1)^2\, e^{\pi i r z} \, \dd z, \label{eqn:J-in-psiS-form}\\
    J_2(r) &= \int_{-1 + i}^{i} \psi_S\!\left(\frac{-1}{z + 1}\right) (z + 1)^2\, e^{\pi i r z} \, \dd z,
    & J_4(r) &= \int_{1 + i}^{i} \psi_S\!\left(\frac{-1}{z - 1}\right) (z - 1)^2\, e^{\pi i r z} \, \dd z, \notag\\
    J_5(r) &= -2\int_{0}^{i} \psi_S\!\left(\frac{-1}{z}\right) z^2\, e^{\pi i r z} \, \dd z. \notag
\end{align}
The proofs of the analytic properties of $b$ below follow those for $a$ with these expressions in place of the $\phi_0$-integrands; we record the statements and indicate only the differences.

\subsection{Schwartzness of $b$}

As in the case of $a(x)$, we start with exponential bounds for the $\psi$-functions.

\begin{lemma}\label{lemma:psi-bound}\uses{lemma:mod-div-disc-bound, lemma:psi-new, lemma:poly-growth-mul, prop:H2-fourier, prop:H3-fourier, prop:H4-fourier}
    There exist constants $C_I, C_S, C_T > 0$ such that
    \begin{align}
        |\psi_I(z)| &\le C_I e^{2\pi \Im z}, \label{eqn:psiI-bound} \\
        |\psi_T(z)| &\le C_T e^{2\pi \Im z}, \label{eqn:psiT-bound} \\
        |\psi_S(z)| &\le C_S e^{- \pi \Im z} \label{eqn:psiS-bound}
    \end{align}
    for all $z$ with $\Im z > 1/2$.
\end{lemma}
\begin{proof}
    We apply Lemma~\ref{lemma:mod-div-disc-bound} to the numerators in Lemma~\ref{lemma:psi-new}. The $H$-functions have Fourier expansions with polynomially growing coefficients (Propositions~\ref{prop:H2-fourier}--\ref{prop:H4-fourier}), and, in the convention of Lemma~\ref{lemma:mod-div-disc-bound}, $n_0(H_2) = 1$ (as $H_2 = 16 q^{1/2} + O(q^{3/2})$) while $n_0(H_3) = n_0(H_4) = 0$. By Lemma~\ref{lemma:poly-growth-mul}, the numerators of $\psi_I$, $\psi_T$ and $\psi_S$ have $n_0 = 0$, $0$ and $3$ respectively, which gives the three bounds.
\end{proof}

\begin{lemma}\label{lemma:bound-J1-J3-J5}\uses{def:b-definition, lemma:psi-bound, lemma:psi-slash-relations}
There exists a constant $C > 0$ such that for all $r \geq 0$,
\begin{align}
    |J_1(r)|, |J_3(r)|, |J_5(r)| &\le C \int_1^{\infty} e^{-\pi s}\, e^{-\pi r / s}\, \dd s.
\end{align}
\end{lemma}
\begin{proof}
Identical to the proof of Lemma~\ref{lem:bound-I1-I3-I5}, starting from the expressions \eqref{eqn:J-in-psiS-form}, with $\phi_0$ replaced by $\psi_S$ and the bound $C_0 e^{-2\pi s}$ of Corollary~\ref{cor:phi0-bound} replaced by the bound $C_S e^{-\pi s}$ of \eqref{eqn:psiS-bound}.
\end{proof}

\begin{lemma}\label{lemma:bound-J2-J4-J6}\uses{def:b-definition, lemma:psi-bound, lemma:psi-slash-relations}
    There exist $C_1, C_2 > 0$ such that for all $r \geq 0$,
    \begin{align}
        |J_2(r)|, |J_4(r)| &\le C_1 e^{-\pi r}, \\
        |J_6(r)| &\le C_2 \frac{e^{-\pi (r + 1)}}{r + 1}.
    \end{align}
\end{lemma}
\begin{proof}
For $J_2$ and $J_4$, argue as for $I_2$ and $I_4$ in Lemma~\ref{lem:bound-I2-I4-I6}: the integrands are continuous, hence bounded, on the compact horizontal contours at height $1$, and $|e^{\pi i r z}| = e^{-\pi r}$ there. For $J_6$, parametrise $z = it$, $t \in [1, \infty)$, and use \eqref{eqn:psiS-bound}:
\[
    |J_6(r)| \le 2 \int_1^\infty |\psi_S(it)|\, e^{-\pi r t}\, \dd t \le 2 C_S \int_1^\infty e^{-\pi(r+1)t}\, \dd t = \frac{2C_S}{\pi} \frac{e^{-\pi(r+1)}}{r+1}.
\]
\end{proof}

Note the asymmetry with the $a$-side: since $\psi_S$ decays like $e^{-\pi \Im z}$ rather than $e^{-2\pi \Im z}$, the integral defining $J_6$ converges (and defines a smooth function) only for $r > -1$, whereas $I_6$ converged for $r > -2$. Arguing exactly as in Lemma~\ref{lem:Ij-smooth}, the functions $J_1, \ldots, J_5$ are smooth on $\R$ and $J_6$ is smooth on $(-1, \infty)$, with all derivatives of all $J_j$ decaying rapidly on $[0, \infty)$ by Lemmas~\ref{lemma:bound-J1-J3-J5}, \ref{lemma:bound-J2-J4-J6} and~\ref{lem:integral-bound}. Corollary~\ref{cor:schwartz-of-nonneg-decay} then applies with $\theta = -1$, the cutoff transitioning on $[-1/2, 0]$.

\begin{proposition}\label{prop:b-schwartz}\uses{lemma:psi-bound}\lean{MagicFunction.FourierEigenfunctions.b}\leanok
$b(x)$ is a Schwartz function.
\end{proposition}
\begin{proof}
\uses{lemma:bound-J1-J3-J5, lemma:bound-J2-J4-J6, lem:integral-bound, cor:schwartz-of-nonneg-decay}
As for Proposition~\ref{prop:a-schwartz}: $b_\rad$ is smooth on $(-1, \infty)$ with all derivatives decaying rapidly on $[0, \infty)$, so Corollary~\ref{cor:schwartz-of-nonneg-decay} with $\theta = -1$ shows that $b = b_\rad(\|\cdot\|^2)$ is Schwartz.
\end{proof}

\subsection{The Eigenfunction Property of $b$}

The Fourier transform also permutes the summands of $b$, but with a sign change:
\begin{align}
    \mathcal{F}(J_1 + J_2) &= -(J_3 + J_4), \label{eqn:b-perm-12}\\
    \mathcal{F}(J_5) &= -J_6. \label{eqn:b-perm-56}
\end{align}

\begin{lemma}\label{lemma:perm-J12}\uses{def:b-definition, lemma:Gaussian-Fourier, lemma:psi-slash-relations, prop:b-schwartz}\lean{MagicFunction.b.Fourier.perm_J₁_J₂}\leanok
\eqref{eqn:b-perm-12} holds: $\mathcal{F}(J_1 + J_2) = -(J_3 + J_4)$.
\end{lemma}
\begin{proof}
As in Lemma~\ref{lemma:perm-I12}, we exchange the Fourier integral with the contour integral (absolute convergence follows from the bounds above), apply Lemma~\ref{lemma:Gaussian-Fourier}, and substitute $w = -1/z$. Here the relevant transformation law is \eqref{eqn:psi-rel-TS}: $\psi_T(-1/w)\, w^2 = -\psi_T(w)$, so the integrand becomes $-\psi_T(w)\, e^{\pi i \|\xi\|^2 w}$, and the same contour deformation as in Lemma~\ref{lemma:perm-I12} identifies the result with $-(J_3 + J_4)$.
\end{proof}

\begin{lemma}\label{lemma:perm-J56}\uses{def:b-definition, lemma:Gaussian-Fourier, lemma:psi-slash-relations, prop:b-schwartz}\lean{MagicFunction.b.Fourier.perm_J₅}\leanok
\eqref{eqn:b-perm-56} holds: $\mathcal{F}(J_5) = -J_6$.
\end{lemma}
\begin{proof}
As in Lemma~\ref{lemma:perm-I56}, using the definition \eqref{eqn:psiS-define} in the form $\psi_I(-1/w)\, w^2 = \psi_S(w)$. The substitution $w = -1/z$ maps the contour of $J_5$ onto that of $J_6$ with reversed orientation, whence the sign.
\end{proof}

\begin{lemma}\label{lemma:perm-J-reverse}\uses{lemma:perm-J12, lemma:perm-J56, lemma:radial-fourier-inversion}\lean{MagicFunction.b.Fourier.perm_₃_J₄, MagicFunction.b.Fourier.perm_J₆}\leanok
$\mathcal{F}(J_3 + J_4) = -(J_1 + J_2)$ and $\mathcal{F}(J_6) = -J_5$.
\end{lemma}
\begin{proof}\leanok
Apply $\mathcal{F}$ to \eqref{eqn:b-perm-12} and \eqref{eqn:b-perm-56} and use $\mathcal{F}^2 = \mathrm{id}$ on radial Schwartz functions (Lemma~\ref{lemma:radial-fourier-inversion}).
\end{proof}

\begin{proposition}\label{prop:b-fourier}\uses{def:b-definition, prop:b-schwartz, lemma:perm-J12, lemma:perm-J56, lemma:perm-J-reverse}\lean{MagicFunction.b.Fourier.eig_b}\leanok
$b(x)$ satisfies \eqref{eqn:b-fourier}.
\end{proposition}
\begin{proof}\leanok
By linearity, Lemmas~\ref{lemma:perm-J12}, \ref{lemma:perm-J56} and \ref{lemma:perm-J-reverse} give
\[
    \mathcal{F}(b) = -(J_3 + J_4) - (J_1 + J_2) - J_6 - J_5 = -b.
\]
\end{proof}

\subsection{Double Zeroes and Special Values of $b$}

Now we regard the radial function $b$ as a function on $\R_{\geq0}$ and check that $b$ has double roots at $\Lambda_8$-points.

\begin{corollary}\label{cor:psiI-near-0-infty}\uses{lemma:psi-bound, lemma:psi-slash-relations}
We have
\begin{align}
    \psi_I(it) &= O(t^2 e^{\pi/t}) \quad \text{as } t \to 0 \label{eqn:psiI-near-0} \\
    \psi_I(it) &= O(e^{2 \pi t}) \quad \text{as } t \to \infty. \label{eqn:psiI-near-infty}
\end{align}
\end{corollary}
\begin{proof}
By \eqref{eqn:psi-rel-SS}, we have
\begin{equation}
    \psi_I(it) = (it)^{2} \psi_S\left(\frac{-1}{it}\right) = -t^{2} \psi_S\left(\frac{i}{t}\right),
\end{equation}
and combined with \eqref{eqn:psiS-bound} we get \eqref{eqn:psiI-near-0}.
\eqref{eqn:psiI-near-infty} follows from Lemma \ref{lemma:psi-bound}.
\end{proof}

We will also need the following identity.

\begin{lemma}\label{lemma:psi-sum}\uses{def:psiI-psiT-psiS}
$\psi_T+\psi_S=\psi_I$.
\end{lemma}
\begin{proof}
From definitions \eqref{eqn:psiI-define}--\eqref{eqn:psiS-define} we get
\begin{align}\psi_T+\psi_S=&(h-h|_{-2}ST)|_{-2}T+(h-h|_{-2}ST)|_{-2}S\notag\\
=&h|_{-2}T-h|_{-2}ST^2+h|_{-2}S-h|_{-2}STS.\notag\end{align}
Note that $ST^2S$ belongs to $\Gamma_0(2)$. Thus, since $h\in M^!_{-2}(\Gamma_0(2))$ we get
$$\psi_T+\psi_S=h|_{-2}T-h|_{-2}STS. $$
Now we observe that $T$ and $STS(ST)^{-1}$ are also in $\Gamma_0(2)$. Therefore,
$$\psi_T+\psi_S=h|_{-2}T-h|_{-2}STS=h-h|_{-2}ST=\psi_I.$$
\end{proof}

\begin{proposition}\label{prop:b-double-zeros}\uses{lemma:psiI-psiT-psiS-fourier, def:psiI-psiT-psiS, cor:psiI-near-0-infty, cor:disc-nonvanishing, lem:CG-unbounded, lemma:psi-sum, lemma:psi-slash-relations}
For $r>\sqrt{2}$ function $b(r)$ can be expressed as
\begin{equation}\label{eqn: b double zeroes}
    b(r)=-4\sin(\pi r^2/2)^2\,\int\limits_{0}^{i\infty}\psi_I(z)\,e^{\pi i r^2 \,z}\,dz.
\end{equation}
\end{proposition}
\begin{proof}
We denote the right hand side of~\eqref{eqn: b double zeroes} by $c(r)$.
By Corollary \ref{cor:psiI-near-0-infty}, the integral in~\eqref{eqn: b double zeroes} converges for $r>\sqrt{2}$.
As in Proposition~\ref{prop:a-double-zeros}, we write $-4\sin(\pi r^2/2)^2 = e^{\pi i r^2} - 2 + e^{-\pi i r^2}$ and shift variables:
$$c(r)=\int\limits_{-1}^{i\infty-1}\psi_I(z+1)\,e^{\pi i r^2 \,z}\,dz-2\int\limits_{0}^{i\infty}\psi_I(z)\,e^{\pi i r^2 \,z}\,dz+
\int\limits_{1}^{i\infty+1}\psi_I(z-1)\,e^{\pi i r^2 \,z}\,dz.$$
By \eqref{eqn:psiT-define} and \eqref{eqn:psi-rel-TT}, $\psi_I(z+1) = \psi_I(z - 1) = \psi_T(z)$.
From the Fourier expansion~\eqref{eqn: psi fourier I} we know that $\psi_I(z)=e^{-2\pi i z}+O(1)$ as $\Im(z)\to\infty$.
By assumption $r^2>2$, hence the integrands tend to $0$ as $\Im(z) \to \infty$ and we can deform the paths of integration by Lemma~\ref{lem:CG-unbounded}:
\begin{align}
\int\limits_{-1}^{i\infty-1}\psi_I(z+1)\,e^{\pi i r^2 \,z}\,dz=&
\int\limits_{-1}^{i}\psi_T(z)\,e^{\pi i r^2 \,z}\,dz+\int\limits_{i}^{i\infty}\psi_T(z)\,e^{\pi i r^2 \,z}\,dz\label{eqn: inside proof 1}\\
\int\limits_{1}^{i\infty+1}\psi_I(z-1)\,e^{\pi i r^2 \,z}\,dz=&
\int\limits_{1}^{i}\psi_T(z)\,e^{\pi i r^2 \,z}\,dz+\int\limits_{i}^{i\infty}\psi_T(z)\,e^{\pi i r^2 \,z}\,dz,\notag
\end{align}
where the first integrals on the right run along the rectangular contours through $-1 + i$ resp.\ $1 + i$. We obtain
\begin{align}\label{eqn: c1}c(r)=&\int\limits_{-1}^{i}\psi_T(z)\,e^{\pi i r^2 \,z}\,dz+\int\limits_{1}^{i}\psi_T(z)\,e^{\pi i r^2 \,z}\,dz
-2\int\limits_{0}^{i}\psi_I(z)\,e^{\pi i r^2 \,z}\,dz\\
&+2\int\limits_{i}^{i\infty}(\psi_T(z)-\psi_I(z))\,e^{\pi i r^2 \,z}\,dz.\nonumber
    \end{align}
By Lemma~\ref{lemma:psi-sum}, $\psi_T - \psi_I = -\psi_S$, so the last integral is $J_6(r^2)$, and
\begin{align}c(r)=&\int\limits_{-1}^{i}\psi_T(z)\,e^{\pi i r^2 \,z}\,dz+\int\limits_{1}^{i}\psi_T(z)\,e^{\pi i r^2 \,z}\,dz
-2\int\limits_{0}^{i}\psi_I(z)\,e^{\pi i r^2 \,z}\,dz\notag\\
&-2\int\limits_{i}^{i\infty}\psi_S(z)\,e^{\pi i r^2 \,z}\,dz\notag\\
=&b(r).\notag
    \end{align}
\end{proof}
At the end of this section we find another integral representation of $b(r)$ for $r\in\R_{\geq0}$ and compute special values of $b$.

\begin{proposition}\label{prop:b-another-integral}\uses{prop:b-double-zeros, lemma:psiI-psiT-psiS-fourier, def:b-definition}
For $r\geq0$ we have
\begin{equation}\label{eqn:b-another-integral}b(r)=4i\,\sin(\pi r^2/2)^2\,\left(\frac{144}{\pi\,r^2}+\frac{1}{\pi\,(r^2-2)}+\int\limits_0^\infty\,\left(\psi_I(it)-144-e^{2\pi t}\right)\,e^{-\pi r^2 t}\,dt\right).\end{equation}
The integral converges absolutely for all $r\in\R_{\geq 0}$.
\end{proposition}
\begin{proof}
The proof is analogous to the proof of Proposition~\ref{prop:a-another-integral}.
First, suppose that $r>\sqrt{2}$. Then by Proposition~\ref{prop:b-double-zeros}
$$b(r)=4i\,\sin(\pi r^2/2)^2\,\int\limits_{0}^{\infty}\psi_I(it)\,e^{-\pi r^2 t}\,dt. $$
From \eqref{eqn: psi fourier I} we obtain
\begin{equation}\label{eqn: psi asymptotic}
\psi_I(it)=e^{2\pi t}+144+O(e^{-\pi t})\quad\mbox{as}\;t\to\infty.
\end{equation}
For $r>\sqrt{2}$ we have
\begin{equation}
\int\limits_0^\infty \left(e^{2\pi t}+144\right)\,e^{-\pi r^2 t}\,dt
=\frac{1}{\pi\,(r^2-2)}+\frac{144}{\pi\,r^2}.\end{equation}
Therefore, the identity \eqref{eqn:b-another-integral} holds for $r>\sqrt{2}$.

On the other hand, from the definition \eqref{eqn:b-definition} we see that $b(r)$ is analytic in some neighborhood of $[0,\infty)$. The asymptotic expansion \eqref{eqn: psi asymptotic} implies that the right hand side of \eqref{eqn:b-another-integral} is also analytic in some neighborhood of $[0,\infty)$. Hence, by the identity principle for analytic functions, the identity \eqref{eqn:b-another-integral} holds on the whole interval $[0,\infty)$. This finishes the proof of the proposition.
\end{proof}
We see from \eqref{eqn:b-another-integral} that $b(r)\in i\R$ for all $r\in\R_{\geq0}$.
Another immediate corollary of this proposition is

\begin{proposition}\label{prop:b0}\uses{prop:b-another-integral}\lean{MagicFunction.b.SpecialValues.b_zero}\leanok
We have $b(0) = 0$.
\end{proposition}
\begin{proof}
As $r \to 0$, the factor $\sin(\pi r^2/2)^2$ vanishes to order $r^4$, while the expression in parentheses in \eqref{eqn:b-another-integral} has only a simple pole in $r^2$ (the term $\frac{144}{\pi r^2}$). Hence the product tends to $0$.
\end{proof}
